\documentclass[10pt]{article}
\usepackage[a4paper,margin=15mm,landscape]{geometry}
\usepackage[T1]{fontenc}
\usepackage{lmodern}
\usepackage{booktabs}
\usepackage{longtable}
\usepackage{array}
\usepackage{enumitem}
\usepackage{microtype}
\setlength{\parindent}{0pt}
\pagestyle{empty}

\begin{document}

{\large\textbf{Supplementary Table 1}\par}\vspace{2pt}
{\textbf{Gene-level evidence for the programmes in Figure 5a.}\par}\vspace{6pt}

{\footnotesize Every gene drawn in Figure 5a, that is, every gene detected in at least 5\%
of the cells its own programme is scored in. \textbf{Detection} is the fraction of those
cells with a non-zero count. \textbf{Principal cell type} is the cell type with the highest
detection of that gene \emph{in this atlas}, computed over the programme's own cell subset
and restricted to types with at least 200 cells; the value in parentheses is detection
within that type. The function column gives one indicative reference per claim and is not an
exhaustive citation. Genes below the detection floor, and the innate-sensing and
cytokine-tone programmes that fail the 60\% measurability gate, are excluded from Figure 5
and from this table.\par}

\vspace{8pt}
{\footnotesize
\begin{longtable}{@{}l r l p{85mm} l@{}}
\toprule
\textbf{Gene} & \textbf{Detection} & \textbf{Principal cell type (detection)} &
\textbf{Function relevant to the programme} & \textbf{Ref.} \\
\midrule
\endhead
\multicolumn{5}{l}{\textbf{Mucus permeability} \;---\; 13 genes scored, empirical $P$ = 0.066, in 92,245 cells} \\[1pt]
\midrule
\textit{MUC16} & 0.74 & Club (0.93) & Largest membrane-tethered mucin; inner, cell-associated barrier layer & [1,2] \\
\textit{MUC4} & 0.74 & Goblet (0.91) & Membrane-tethered mucin of the periciliary layer & [1,2] \\
\textit{MUC1} & 0.71 & Goblet inflammatory (0.91) & Membrane-tethered mucin of the periciliary brush; also limits and resolves airway inflammation during infection & [1,2] \\
\textit{CLDN4} & 0.56 & Club (0.72) & Tight-junction claudin of the apical junctional complex & [4,5] \\
\textit{MUC5AC} & 0.52 & Goblet (0.82) & Gel-forming secreted mucin; concentrated in proximal airways and made by CCSP-positive secretory and goblet cells & [1,3] \\
\textit{TJP1} & 0.37 & Goblet inflammatory (0.53) & ZO-1; scaffold linking tight-junction strands to the actin cytoskeleton & [4,5] \\
\textit{SCGB1A1} & 0.35 & Club (0.70) & Club cell secretory protein CC16; anti-inflammatory and anti-oxidant component of airway secretions & [7] \\
\textit{CLDN1} & 0.29 & Goblet inflammatory (0.49) & Tight-junction claudin; its loss raises paracellular permeability in airway epithelium & [4,5] \\
\textit{LYZ} & 0.22 & Goblet (0.39) & Lysozyme; principal muramidase of airway surface liquid & [6] \\
\textit{BPIFA2} & 0.16 & Goblet inflammatory (0.32) & Secreted BPI-fold protein of the salivary/airway secretome with anticandidal activity & [8] \\
\textit{OCLN} & 0.15 & Goblet inflammatory (0.25) & Occludin; integral tight-junction protein & [4,5] \\
\textit{GALNT4} & 0.10 & Club (0.14) & GalNAc-transferase that initiates mucin-type O-glycosylation, the modification that gives mucins their gel properties & [9] \\
\textit{MUC5B} & 0.06 & Club (0.40) & Gel-forming secreted mucin; dominant mucin of healthy superficial epithelium and submucosal glands, required for mucociliary clearance & [1,3] \\
\addlinespace
\multicolumn{5}{l}{\textbf{Antigen presentation} \;---\; 11 genes scored, empirical $P$ = 0.000, in 68,179 cells} \\[1pt]
\midrule
\textit{HLA-B} & 0.92 & Goblet (0.96) & MHC class I heavy chain; presents cytosolic peptides to CD8+ T cells & [10] \\
\textit{HLA-C} & 0.92 & Goblet (0.95) & MHC class I heavy chain; presents cytosolic peptides to CD8+ T cells & [10] \\
\textit{HLA-A} & 0.88 & Goblet (0.92) & MHC class I heavy chain; presents cytosolic peptides to CD8+ T cells & [10] \\
\textit{CD74} & 0.87 & Macrophage (0.94) & Invariant chain; chaperones MHC class II and blocks the groove until endosomal peptide loading & [10] \\
\textit{LY6E} & 0.80 & Goblet inflammatory (0.89) & GPI-anchored interferon-stimulated protein that modulates a late step of viral entry & [11] \\
\textit{HLA-DRA} & 0.64 & Macrophage (0.85) & MHC class II alpha chain; presents endosomal peptides to CD4+ T cells & [10] \\
\textit{HLA-DRB1} & 0.58 & Monocyte (0.76) & MHC class II beta chain; presents endosomal peptides to CD4+ T cells & [10] \\
\textit{IFNGR1} & 0.32 & Goblet inflammatory (0.37) & IFN-gamma receptor ligand-binding chain; complete deficiency causes susceptibility to mycobacterial and other intracellular infection & [14] \\
\textit{IFNGR2} & 0.27 & Monocyte (0.43) & IFN-gamma receptor signalling chain, partner of IFNGR1 & [14] \\
\textit{CD40} & 0.15 & Goblet inflammatory (0.18) & Costimulatory receptor; CD40-CD40L engagement licenses antigen-presenting cells and supports B-cell class switching & [13] \\
\textit{LRP1} & 0.10 & Monocyte (0.38) & CD91; endocytic scavenger receptor that internalises ligand complexes for intracellular processing & [12] \\
\addlinespace
\multicolumn{5}{l}{\textbf{IgA support} \;---\; 3 genes scored, empirical $P$ = 0.002, in 72,770 cells} \\[1pt]
\midrule
\textit{PIGR} & 0.84 & Goblet (0.96) & Polymeric immunoglobulin receptor; transcytoses dimeric IgA to the mucosal surface and is made predominantly by secretory cells & [15] \\
\textit{TNFSF13} & 0.06 & Goblet inflammatory (0.07) & APRIL; drives T-cell-independent IgA class-switch recombination at mucosal sites & [17] \\
\textit{TGFB1} & 0.06 & B mem (0.28) & Switch factor for IgA class-switch recombination; also down-regulates epithelial pIgR and so reduces secretory IgA transport & [16,17] \\
\addlinespace
\bottomrule
\end{longtable}
}

\vspace{4pt}
{\footnotesize\textbf{References}\par}
{\scriptsize
\begin{itemize}[leftmargin=2.2em,itemsep=1pt,topsep=2pt]
\item[1.] Ma J, Rubin BK, Voynow JA. Mucins, mucus, and goblet cells. Chest 154:169-176 (2018). \textsc{doi}: \texttt{10.1016/j.chest.2017.11.008}
\item[2.] Kim KC. Role of epithelial mucins during airway infection. Pulm Pharmacol Ther 25:415-419 (2012). \textsc{doi}: \texttt{10.1016/j.pupt.2011.12.003}
\item[3.] Okuda K et al. Localization of secretory mucins MUC5AC and MUC5B in normal/healthy human airways. Am J Respir Crit Care Med 199:715-727 (2019). \textsc{doi}: \texttt{10.1164/rccm.201804-0734OC}
\item[4.] Gon Y, Hashimoto S. Role of airway epithelial barrier dysfunction in pathogenesis of asthma. Allergol Int 67:12-17 (2018). \textsc{doi}: \texttt{10.1016/j.alit.2017.08.011}
\item[5.] Gan H et al. Protein kinase D promotes airway epithelial barrier dysfunction and permeability through down-regulation of claudin-1. J Biol Chem 288:37343-37354 (2013). \textsc{doi}: \texttt{10.1074/jbc.M113.511527}
\item[6.] Ganz T. Antimicrobial polypeptides. J Leukoc Biol 75:34-38 (2004). \textsc{doi}: \texttt{10.1189/jlb.0403150}
\item[7.] Almuntashiri S et al. Club cell secreted protein CC16: potential applications in prognosis and therapy for pulmonary diseases. J Clin Med 9:4039 (2020). \textsc{doi}: \texttt{10.3390/jcm9124039}
\item[8.] Khovidhunkit W et al. Parotid secretory protein is an HDL-associated protein with anticandidal activity. Am J Physiol Regul Integr Comp Physiol 288:R1306-R1315 (2005). \textsc{doi}: \texttt{10.1152/ajpregu.00007.2004}
\item[9.] Kato K, Hansen L, Clausen H. Polypeptide N-acetylgalactosaminyltransferase-associated phenotypes in mammals. Molecules 26:5504 (2021). \textsc{doi}: \texttt{10.3390/molecules26185504}
\item[10.] Pishesha N, Harmand TJ, Ploegh HL. A guide to antigen processing and presentation. Nat Rev Immunol 22:751-764 (2022). \textsc{doi}: \texttt{10.1038/s41577-022-00707-2}
\item[11.] Mar KB et al. LY6E mediates an evolutionarily conserved enhancement of virus infection by targeting a late entry step. Nat Commun 9:3603 (2018). \textsc{doi}: \texttt{10.1038/s41467-018-06000-y}
\item[12.] Hvidberg V et al. Identification of the receptor scavenging hemopexin-heme complexes [LRP1/CD91]. Blood 106:2572-2579 (2005). \textsc{doi}: \texttt{10.1182/blood-2005-03-1185}
\item[13.] Grewal IS, Flavell RA. The role of CD40 ligand in costimulation and T-cell activation. Immunol Rev 153:85-106 (1996). \textsc{doi}: \texttt{10.1111/j.1600-065x.1996.tb00921.x}
\item[14.] Reuter U et al. Correction of complete interferon-gamma receptor 1 deficiency by bone marrow transplantation. Blood 100:4234-4235 (2002). \textsc{doi}: \texttt{10.1182/blood-2002-02-0433}
\item[15.] Blackburn JB et al. Secretory cells are the primary source of pIgR in small airways. Am J Respir Cell Mol Biol 67:334-345 (2022). \textsc{doi}: \texttt{10.1165/rcmb.2021-0548OC}
\item[16.] Wyatt TA et al. Malondialdehyde-acetaldehyde adduct formation decreases immunoglobulin A transport across airway epithelium in smokers who abuse alcohol. Am J Pathol 191:1732-1742 (2021). \textsc{doi}: \texttt{10.1016/j.ajpath.2021.06.007}
\item[17.] Tezuka H et al. Regulation of IgA production by naturally occurring TNF/iNOS-producing dendritic cells. Nature 448:929-933 (2007). \textsc{doi}: \texttt{10.1038/nature06033}
\end{itemize}
}

\end{document}
